% Actual finite-point restrictions for the S17 norm sheaf.
% GPT-6 Astra Pro; /root/fixed_surface; 22 November 2026 UTC.
\section[HD-F1-S19: norm-fourfold point stabilizers]{HD-F1-S19: the norm sheaf has no tensor stabilizer and every finite point stabilizer has order dividing six}
\label{sec:f1-norm-fourfold-point-stabilizers}

For the actual normal norm image $i:Y\hookrightarrow P$ and
$a:T\to P$ of S13--S18, pullback gives an injective homomorphism
of abelian varieties
\[
                    a^*:\operatorname{Pic}^0(P)
                              \hookrightarrow\operatorname{Pic}^0(T).
 \tag{F1S19.1}
\]
Equivalently, after fixing a base point, the actual Albanese
homomorphism $\operatorname{Alb}(T)\to P$ is surjective with
connected kernel. No dimension for that kernel is asserted here.
There is no nontrivial degree-zero line bundle on $P$ whose
restriction to $Y$ is trivial.

Let $G\subset P\times\widehat P$ be any finite subgroup defining
a faithful point action, and suppose the actual sheaf
$E=i_*\mathcal O_Y$ has a coherent lift for that action. Then
\[
                              |G|\mid6.
 \tag{F1S19.2}
\]
This is a restriction on actual point stabilizers, independent of
semiregularity. The proof does not assert existence of a lift at
any of the possible orders.

\paragraph{1. Pullback on the Picard varieties and the actual cyclic cover.}
The differential of (F1S19.1) is
$H^1(P,\mathcal O_P)\to H^1(T,\mathcal O_T)$, which is injective
by S18. Its kernel is therefore finite. Every member of this
kernel is a torsion line bundle.

Write $a=ib$. The identity $Rb_*\mathcal O_T=\mathcal O_Y$
from S17 and the projection formula give, for every line bundle
$L$ on $P$,
\[
 Rb_*a^*L=(i^*L)\otimes^{\mathbf L}Rb_*\mathcal O_T=i^*L.
 \tag{F1S19.3}
\]
Consequently $a^*L\simeq\mathcal O_T$ implies
$i^*L\simeq\mathcal O_Y$; the reverse implication follows by
pullback. Thus the two actual restriction kernels agree.

Suppose a nontrivial element $L$ of this finite kernel has exact
order $r>1$. The subgroup $\langle L\rangle\subset\widehat P$
defines the quotient isogeny
$\widehat P\to\widehat P/\langle L\rangle$. Its dual is the
actual cyclic etale isogeny
\[
             q:P'=\bigl(\widehat P/\langle L\rangle\bigr)^\wedge
                          \longrightarrow P,\qquad\deg q=r.
 \tag{F1S19.4}
\]
The eigensheaf decomposition of its finite pushforward is
$q_*\mathcal O_{P'}=\bigoplus_{j=0}^{r-1}L^j$, with the
multiplication given by the chosen trivialization of $L^r$.
On $Y$ every $L^j$ is trivial. A trivialization of $L|_Y$ makes
the multiplication constant: its $r$th power differs from the
fixed trivialization of $L^r$ by a nonzero element of
$H^0(Y,\mathcal O_Y)=\mathbb C$. Rescale it by an $r$th root
of that constant. The resulting algebra is the split cyclic
algebra. Thus the actual inverse image is a disjoint union
\[
                      q^{-1}Y=Y'_1\sqcup\cdots\sqcup Y'_r,
                          \qquad Y'_j\simeq Y.
 \tag{F1S19.5}
\]
The deck transformations of (F1S19.4) are translations. They
act transitively on these components and trivially on the
cohomology of $P'$, so all $[Y'_j]$ are equal. Since the
components are disjoint and have half the dimension of $P'$,
\[
                  \int_{P'}[Y'_1]^2
                     =\int_{P'}[Y'_1][Y'_2]=0.
\]
On the other hand, $q^*[Y]=r[Y'_1]$ and the finite-degree
projection formula give
\[
   r\cdot648
       =r\int_P[Y]^2
       =\int_{P'}(q^*[Y])^2
       =r^2\int_{P'}[Y'_1]^2=0,
 \tag{F1S19.6}
\]
contrary to the actual self-intersection $648$ proved in S14.
This proves (F1S19.1) and the asserted tensor-stabilizer vanishing.

The dual Albanese homomorphism is surjective because its dual
has injective differential. To check connectedness without
discarding a possible isogeny, divide $\operatorname{Alb}(T)$
by the identity component of that kernel. The resulting map
to $P$ is an isogeny. The kernel of its dual is a subgroup of
the kernel in (F1S19.1), hence is zero. This isogeny has degree
one. The original Albanese kernel is therefore connected, as
claimed.

\paragraph{2. Every finite point action injects into translations.}
For $(x,L)\in G$, the actual lift gives an isomorphism
\[
                    t_x^*i_*\mathcal O_Y\otimes L
                                      \simeq i_*\mathcal O_Y.
 \tag{F1S19.7}
\]
Equality of supports shows that $x$ stabilizes $Y$ by translation.
If $x=0$, full faithfulness of the pushforward of the closed
immersion identifies (F1S19.7) with
$L|_Y\simeq\mathcal O_Y$. Paragraph 1 forces $L$ to be trivial.
Faithfulness of the specified point subgroup therefore makes
the projection
\[
                  G\longrightarrow P,\qquad (x,L)\longmapsto x
 \tag{F1S19.8}
\]
injective. Write its image as $H$, so $|H|=|G|$. This step
retains the given lift and its coherence; it replaces neither
the functors nor their linearizations.

\paragraph{3. The two exact divisibilities.}
Take the actual quotient isogeny $p:P\to P/H$ and let
$Y_0=p(Y)$ with its reduced structure. Because $H$ stabilizes
$Y$, and $p$ is etale, one has equality of reduced schemes
$p^{-1}Y_0=Y$. Thus $Y\to Y_0$ is finite etale of degree
$h=|H|$. Flat pullback of its fundamental cycle is $[Y]$, and
\[
          648=\int_P[Y]^2
                     =h\int_{P/H}[Y_0]^2.
 \tag{F1S19.9}
\]
The last self-intersection is an integer. Hence $h\mid648$.

We also justify holomorphic-Euler multiplicativity on these
actual singular varieties. Let $T_0\to Y_0$ be a projective
resolution. The base change $T'=T_0\times_{Y_0}Y$ is smooth,
finite etale of degree $h$ over $T_0$, and proper birational
over $Y$. Thus it is a resolution of the normal variety $Y$.
Rational singularities of $Y$ give
$R(T'\to Y)_*\mathcal O_{T'}=\mathcal O_Y$. Proper flat
base change along the faithfully flat map $Y\to Y_0$ then
gives $R(T_0\to Y_0)_*\mathcal O_{T_0}=\mathcal O_{Y_0}$.
In particular both resolutions preserve holomorphic Euler
characteristic. On the smooth etale cover $T'\to T_0$, the
tangent bundle is the pullback of $T_{T_0}$, so the actual
Riemann--Roch and degree formulas give
\[
 \chi(\mathcal O_Y)=\chi(\mathcal O_{T'})
     =h\chi(\mathcal O_{T_0})=h\chi(\mathcal O_{Y_0}).
 \tag{F1S19.10}
\]
S13 and S17 give $\chi(\mathcal O_Y)=318$, so $h\mid318$.
Combining (F1S19.9)--(F1S19.10) gives
$h\mid\gcd(648,318)=6$, proving (F1S19.2).

\paragraph{Framework and scope.}
All maps are the displayed Picard dual isogeny, its actual
eigensheaf algebra, the quotient by the projected translation
group, the resolution base-change square, and the actual
coherent point-linearization isomorphisms. The scalar
characteristic maps used to prove injectivity on tangent spaces
are those of S18. Flat pullback and proper pushforward satisfy
the same $K_0$, GRR, Poincare and Adams identities already
retained in S13--S18. The determinant class is not changed in
the contradiction with (F1S19.6); its fixed numerical input
is the computed fundamental-cycle square of this specified $Y$.
The order bound does not supply a full trace calculation or a
semiregularity verdict for $\mathcal O_Y$.
